\documentclass[main]{subfiles}
\usepackage[sols]{workbook}
\numbering{ws3-5}

\begin{document}

\ifSubfilesClassLoaded{%
  \chaptermark{\chapterthree}
}{}

\section{Mixing Problems and Intro to Second-Order Differential Equations} \label{ws3-5}

\begin{framed}
\centerline{\textbf{Lesson Goals}}

You should be able to:
\begin{itemize}
\item Use first-order differential equations to model and analyze various mixing problems
\item Determine when qualitative analysis can and cannot be used to analyze mixing problems
\item Use qualitative analysis when appropriate to determine the long-term behavior of a mixed system.
\item Determine the appropriate amount initial conditions necessary for a second order differential equation and solve a differential equation involving free fall.
\end{itemize}
\end{framed}

\begin{enumerate}
    \item
  \note{A good thing to discuss is what
  should happen in the long run; that is, what do they think
  $\displaystyle \lim_{t \to \infty} S(t)$ is?}
  
  You are mixing up a big batch
  of eggnog for the holidays.  You have a big mixing vat
  containing 100 liters of egg yolks.  At time $t = 0$, you
  start pouring a milk and sugar mixture containing 80 grams of sugar per
  liter into the vat at a rate of 4 liters per minute.  At the same  
  time, mixed eggnog starts leaving the vat a rate of 4 liters per  
  minute.  The fluid in the vat is continuously stirred and well-mixed,
  so that the sugar concentration remains uniform throughout the vat.

  \begin{enumerate}
  \item
  \begin{problem}[2.5in]
  Write a differential equation for the amount $S(t)$ of sugar (measured
  in grams) in the vat at time $t$ including an initial condition.  Can you tell what happens with $S(t)$ in the long run?
  \end{problem}
  
  \begin{solution}
    It is easy to overthink this problem, but here is some guidance that
    will help prevent you from going astray.  Remember that the rate of 
    change of the amount of sugar in the vat equals the rate at which sugar 
    enters the vat minus the rate at which sugar leaves the vat:
    $$
      \frac{dS}{dt} \; = \; \mbox{rate in of sugar} - \mbox{rate out of sugar}.
    $$
    The problem statement gives us enough information to compute the rate in
    and the rate out, and we'll handle those rates separately.

    Regarding the rate in:  We know the rate at which fluid enters the vat, 
    and also the sugar concentration of the fluid entering the vat.  We
    are told that fluid enters the vat at 4 liters per minute, and that the
    fluid contains 80 grams per liter of sugar.  This means that the rate at
    which sugar enters the vat is:
    $$
      \mbox{rate in} =  \left( 4 \: \frac{\mbox{liters}}{\mbox{minute}} \right) \cdot
      \left( 80 \: \frac{\mbox{grams}}{\mbox{liter}} \right) \; = \; 320 \:
      \frac{\mbox{grams}}{\mbox{minute}}.
    $$

    Regarding the rate out:  We know the rate at which fluid leaves the vat,
    namely 4 liters per minute, is the same as the rate the fluid enters the vat.
    This means the volume of fluid in the vat remains constant at 100 liters.
    In order to know the rate at which sugar leaves
    the vat, we need to know the sugar concentration of what's leaving the vat.
    We know that  $S(t)$ represents the amount of sugar in the vat, and that the
    vat always contains 100 liters of fluid, and that the fluid in the vat is
    well-mixed by the continuous stirring.  This means that the fluid leaving the
    vat has a sugar concentration of $S(t)/100$, that is, the amount of sugar
    divided by the volume of fluid in the vat.  This means that the rate at which
    sugar leaves the vat is given by
    $$
      \mbox{rate out} =  \left( 4 \: \frac{\mbox{liters}}{\mbox{minute}} \right) \cdot
      \left( \frac{S(t)}{100} \: \frac{\mbox{grams}}{\mbox{liter}} \right) \; = \; \frac{4 S(t)}{100} \:
      \frac{\mbox{grams}}{\mbox{minute}}.
    $$

    Putting this all together, we have
    $$
      \frac{dS}{dt} = 320 - \frac{4 S}{100}.
    $$ 
    The initial condition is \fbox{$S(0) = 0$} because initially there is no
    sugar in the vat.  

    {\bf Remark:}  This differential equation is separable.  If you separate
    the variables and use the initial condition $S(0) = 0$, you will find a
    formula for $S(t)$, the amount of sugar in the vat at time $t$.  In case
    you're interested, the solution is $\displaystyle S(t) = 8000 \left( 1 - e^{-t/25} \right)$.
    As $t \to \infty$, note that $S(t) \to 8000$.  The equilibrium $S = 8000$ is stable,
    and it makes sense for the following reason:  in the long run, we expect
    the concentration of sugar in the fluid leaving the vat to match the
    concentration in the fluid entering the vat.  If the concentration entering
    is 80 grams/liter and the vat always contains 100 liters of fluid, then there
    should be $80 \cdot 100 = 8000$ grams of sugar in the vat at equilibrium.
  \end{solution}
  
  \item
  \begin{problem}[1.5in]
  Suppose instead that, in the previous problem, mixed eggnog left the  
  vat at a rate of 6 liters per minute.  How would this change the
  differential equation and initial condition?
  \end{problem}
  
  \begin{solution}
  The basic framework is still the same:
  \[ (\textrm{rate of change of sugar in vat}) = (\textrm{rate at which
  sugar enters vat}) - (\textrm{rate at which sugar leaves vat}).\]
  The rate at which sugar enters the vat is exactly the same as before,
  so it is still 320 grams per minute.  The rate at which sugar leaves the
  vat can still be calculated as
  \begin{align*}
  {} & (\textrm{rate at which sugar is leaving vat at time $t$}) \\
  = {} & (\textrm{rate at which eggnog is leaving vat at time $t$})
  \cdot (\textrm{concentration of sugar in eggnog at time $t$}) \\
  = {} & 6 \cdot \left( \frac{\textrm{amount of sugar in
  vat at time $t$}}{\textrm{volume of eggnog in vat at time
  $t$}} \right) \\
  = {} & 6 \cdot \left( \frac{S}{\textrm{volume of eggnog in vat at
  time $t$}}
  \right)
  \end{align*}
  Since stuff is being added to the vat at 4 liters per minute, but
  stuff is leaving the vat at 6 liters per minute, the volume of eggnog
  in the vat is being reduced by 2 liters per minute.  Thus, after $t$
  minutes, the volume remaining in the vat is $100 - 2t$ liters.  Therefore, the
  rate at which sugar is leaving the vat at time $t$ is $6 \left(
  \frac{S}{100 - 2t} \right)$.  So, we have the final differential
  equation $\boxed{\frac{dS}{dt} = 320 - 6 \frac{S}{100 - 2t}}$.  The
  initial condition is still the same, $\boxed{S(0) = 0}$ since there is
  initially no sugar in the vat.

  {\bf Remark:}  The vat is empty when $t = 50$ minutes, after which point
  the differential equation fails as a model.
  \end{solution}
  \end{enumerate}

\pspace{\newpage}

\item The Swedish Chef  (a character on The Muppet Show) is mixing a big pot of turkey-duck-chicken stock to
    make into gravy for his Thanksgiving turducken.  He starts with a
    pot filled with $20$ quarts of liquid, consisting of $10\%$ turkey
    stock and $90\%$ duck stock.  His assistant pours in a
    chicken-turkey stock blend ($40\%$ chicken stock and $60\%$ turkey
    stock) at a rate of $4$ quarts per hour while the Swedish Chef stirs
    vigorously.  Being extremely absent-minded, the Swedish Chef fails
    to notice that he's splashing stock out of the pot at a rate of $4$
    quarts per hour.

\begin{enumerate}
\item \label{swedish-chef-constant-volume}
\begin{problem}[\vfill]
      Write a differential equation whose solution
    $T(t)$ is the amount of turkey stock in the pot at time $t$ hours, where $t = 0$ is the time when the chef's assistant starts pouring and
    the Swedish Chef starts stirring.  Also, write an appropriate initial
    condition.
    \end{problem}

    \begin{solution}
    Since $T(t)$ is the amount of turkey stock in the pot at time $t$,
    $\frac{dT}{dt}$ is the rate of change of the amount of turkey stock in
    the pot.  We know that this is equal to (rate at which turkey stock is
    entering the pot) minus (rate at which turkey stock is leaving the
    pot).

    Let's first focus on the stuff entering the pot.  This is a
    chicken-turkey blend, entering at a rate of 4 quarts per hour.  Only
    $60\%$ of this is turkey stock though, so turkey stock is entering the
    pot at a rate of $(4)(.6) = 2.4$ quarts per hour.  (Here, we are
    using the basic principle (rate of turkey stock entering) = (rate of
    mixture entering) times (concentration of turkey stock in the
    mixture).)

    Now, let's look at the stuff exiting the pot.  We have a similar
    formula: (rate of turkey stock leaving) = (rate of mixture leaving)
    times (concentration of turkey stock in the mixture).  We know that
    the mixture is leaving at a rate of 4 quarts per minute, so we just
    need to figure out the concentration of turkey stock in the mixture.
    This is simple if we think about what our variables mean: $T(t)$ is
    the amount of turkey stock in the pot at time $t$, while the pot
    always contains 20 quarts of stock.  So, the concentration of turkey
    stock in the mixture is $\frac{T(t)}{20}$.  Therefore, turkey stock is
    leaving the pot at a rate of $4 \cdot \frac{T(t)}{20} =
    \frac{T(t)}{5}$ quarts per hour.

    So, our final differential equation is \fbox{$\frac{dT}{dt} = 2.4 -
    \frac{T}{5}$}.  We also have an initial condition, since we know how
    much turkey stock is in the pot at the beginning: $10\%$ of the $20$
    quarts, or 2 quarts.  So, our initial condition is \fbox{$T(0) = 2$}.
    \end{solution}

    \item
    \begin{problem}[\vfill]
    Use qualitative analysis to sketch $T(t)$. 
    \end{problem}

    \begin{solution}
    The following picture shows a family of representative solutions in
    black, together with the graph of $T(t)$ in red.  To find $T(t)$,
    we just use the initial condition $T(0) = 2$.
    \begin{center}
    \includegraphics[width=2in]{ws/images/swedish-chef-constant-volume-qa} 
    \end{center}
    \end{solution}

    \item
    \begin{problem}[\vfill]
    Based on your sketch, what happens to the amount of turkey stock in
    the pot as time goes on? Why does this result make sense?
    \end{problem}

    \begin{solution}
    $T(t)$ is the amount of turkey stock in the pot at time $t$, and we
    can see that this increases and tends to $12$ quarts. This makes sense because eventually the concentration of turkey stock in the pot will match the concentration of turkey stock in the incoming chicken-turkey blend. Since the incoming concentration of turkey stock is 60\%, and the pot always contains 20 quarts of liquid, we should expect that eventually, 60\% of those 20 quarts, or 12 quarts, will be turkey stock.
    \end{solution}
\pspace{\newpage}
    \item \label{swedish-chef-non-constant-volume}
    \begin{problem}[2.25in]
    The Swedish Chef accidentally knocks over his first pot of stock and
    has to start over.  He follows the exact same procedure as before.
    However, being rather agitated over the loss of his first pot of
    stock, he now splashes stock out of the pot at a rate of $5$ quarts
    per hour.  Write a differential equation modeling the new situation.
    %(\textit{Note:} We cannot use qualitative analysis on this differential equation because it is not autonomous; that is, $\frac{dT}{dt}$ depends on $t$.)
    \end{problem}

    \begin{solution}
    The basic framework is exactly the same as in
    \ref{swedish-chef-constant-volume}.  
    We still have $\frac{dT}{dt} = $ (rate at which turkey stock is
    entering the pot) minus (rate at which turkey stock is leaving the
    pot), and the rate at which turkey stock is entering the pot is
    exactly the same as in the previous problem, 2.4 quarts per hour.

    Let's look at the stuff exiting the pot.  Again, we need to use the
    formula (rate at which turkey stock is leaving) = (concentration of
    turkey stock in the mixture) times (rate at which mixture is leaving).
    The rate at which the mixture is leaving is 5 quarts per hour.  The
    concentration of turkey stock in the mixture is equal to (amount of
    turkey stock in mixture) divided by (total amount of mixture).  The
    amount of turkey stock is exactly $T(t)$.  The total amount of mixture
    is a little more complicated.  At time $t = 0$, there is 20 quarts of
    stock in the pot.  However, because the stock is entering at a rate of
    4 quarts per minute and leaving at 5 quarts per minute, there is a net
    loss of 1 quart per minute.  Thus, after $t$ minutes, the amount of
    juice in the pot is $20 - t$.  So, the concentration of turkey stock
    in the mixture at time $t$ is $\frac{T(t)}{20 - t}$, and the rate at
    which turkey stock is leaving the pot is $\frac{T(t)}{20 - t} \cdot
    5$.

    So, our final differential equation is \fbox{$\frac{dT}{dt} = 2.4 -
    \frac{5T}{20 - t}$}.  Again, we have the initial condition $T(0) =
    2$.
    
    Note that we cannot use qualitative analysis on this differential equation, because it is not autonomous.
    \end{solution}

  \end{enumerate}
  
  
  
  
  
  
  
   \item \label{spring2005-final-2b}

    \begin{problem}[3in]
      Due to a malfunctioning mechanism in an industrial plant upriver from
      a lake, the lake has been contaminated with pollutants.  At a time
      we'll designate as $t = 0$, there are $8 \times 10^5$ grams of
      pollutant in the lake.  The malfunction is now being addressed, and the
      concentration of pollutants in the incoming water is decreasing.  The
      concentration of pollutants in the incoming water at time $t$ is given
      by $\displaystyle \frac{1}{100t^2 + 1}$ grams per liter, where $t$ is measured in days.  The $9
      \times 10^7$ liter lake is fed by the river at a rate of 50,000 liters per day, and
      the well-mixed water leaves the lake at the same rate.  
      
      Let $P(t)$ be the
      number of grams of pollutant in the lake at time $t$.  Write a
      differential equation involving $\frac{dP}{dt}$ that models the
      situation.  Include an initial condition. Is it possible to apply qualitative analysis to your differential equation?
          \end{problem}

    \begin{solution}
     We know that:
      \[ \frac{dP}{dt} = (\textrm{rate at which pollutant enters lake}) -
      (\textrm{rate at which pollutant leaves lake}).\]
      \begin{itemize}
      \item rate in of water = 50,000 liters per day      
      \item rate out of water = 50,000 liters per day
      \item concentration of pollutant in incoming water = $\displaystyle
        \frac{1}{100t^2 + 1}$ grams per liter
      \item concentration of pollutant in outgoing water = concentration of
        pollutant in lake, which is just $\displaystyle \frac{\textrm{amount
            of pollutant in lake}}{\textrm{amount of water in lake}} =
        \frac{P}{9 \times 10^7}$.  (The amount of water in the lake is
        always $9 \times 10^7$, since water enters and leaves at the same
        rate.) 
      \end{itemize}
      Therefore,
      \begin{itemize}
      \item rate in of pollutant = (rate in of water) (concentration of
        pollutant in incoming water) = $\displaystyle \frac{50000}{100t^2 +
          1}$ grams per day
      \item rate out of pollutant = (rate out of water) (concentration of
        pollutant in outgoing water) = $\displaystyle 50000 \cdot \frac{P}{9
          \times 10^7} = \frac{5P}{9 \times 10^3}$. 
      \end{itemize}
      Putting this all together, we have the differential equation
      $\boxed{\frac{dP}{dt} = \frac{50000}{100t^2 + 1} - \frac{5}{9 \times
          10^3} P}$.  The fact that there are initially $8 \times 10^5$
      grams of pollutant gives the initial condition $\boxed{P(0) = 8 \times
        10^5}$. 
        
        This differential equation is not autonomous, since $\frac{dP}{dt}$ depends on both $P$ and $t$. Thus, we cannot use qualitative analysis on it.
    \end{solution}

\pspace{\newpage}

\probonly{
 \begin{framed}
  The {\bf order} of a differential equation is the order of the highest derivative that
  appears in the DE.  For instance, the DE $y'' + 6 y' + \cos(y) = t^2$ is of order $2$ (or {\em second-order})
  because the highest derivative is a second derivative.  The DE $\displaystyle ( y' )^{8}
  + 3y = t^{3}$ is of order $1$ (or {\em first-order}) because the highest derivative is a first
  derivative.  
 \end{framed}
}

\item  \label{second-order}

%\note{Students often try to solve partway to write an equation for $s'(t)$ instead of just writing the equation for $s''(t)$.  It's worthwhile to note that we have two initial conditions here and that this has to do with the fact that our differential equation involves a second derivative.}

\note{This is a problem to start discussing the physical intuition about the need for more than one initial condition when we have a second order differential equation. We will be working more thoroughly with 2nd order DEs the following two classes}

  \begin{enumerate}
    \item
    \begin{problem}[1.5in]
    Write a differential equation for $s(t)$ with initial condition(s) that models 
    the following situation:

    A yellow rubber duck is dropped out of the window of an apartment
    building from a height of 80 m.  Let $s = s(t)$ be the height of the
    duck above the ground at time $t$.  (Acceleration due to gravity is 
    $-9.8 \textrm{ m/s}^2$.)  

  \probonly{\textit{Hint:} Write an equation involving a second
    derivative.}
  \end{problem}

  \begin{solution}
    Due to gravity, the rubber duck has a constant acceleration of $-9.8$   
    m/s$^2$.  Therefore, $\boxed{s''(t) = -9.8}$.  We know that the rubber
    duck starts at 80 feet above the ground.  If we let $t = 0$ be the
    time that the rubber duck is dropped, then \fbox{$s(0) = 80$}.  We
    also know that the duck is not moving at the beginning, so its initial
    velocity is 0.  That is, \fbox{$s'(0) = 0$}.
  \end{solution}
    
  \item
  \begin{problem}[2in]
    How many initial conditions did you include in your solution to part (a)?  Why does it make sense that you'd
    need this many to give a complete description of the duck's height over time?
  \end{problem}
  \begin{solution}
    We needed both the initial position $s(0)=80$, and the initial velocity $s'(0)$.
    Without knowing both the initial position and the initial velocity, we would
    have no way of predicting the duck's height above the ground.
  \end{solution}

\item 
\begin{problem}
    Obtain a solution to the differential equation above subject to the given initial conditions.
\end{problem}
\begin{solution}
    Since $s''(t) = -9.8$ we must have $s'(t) = C_1-9.8 t$ and $\displaystyle s(t) = C_2 + C_1 t -\frac{9.8}{2} t^2$. Since $s(0) = C_2$, in order to satisfy $s(0) = 80$ we must have $C_2 = 80$. Similarly, since $s'(0) = C_1$, in order to satisfy $s'(0) = 0$ we must have $C_1 = 0$. Thus, the solution must be $\displaystyle\boxed{s(t)=80 - \frac{9.8}{2}t^2}$
\end{solution}
  


  \end{enumerate}








    
\end{enumerate}


\end{document}